\documentclass[10pt]{article}

\usepackage[utf8]{inputenc}

\usepackage[T1]{fontenc}

\usepackage{amsmath}

\usepackage{amsfonts}

\usepackage{amssymb}

\usepackage[version=4]{mhchem}

\usepackage{stmaryrd}



\begin{document}

жет компенсироваться понижением электронных уровней. Удвоение периода при $\rho=1 / 2$ было впервые обнаружено P. Пайерлсом (R. Peierls). В общем случае $\rho \neq 0,1 / 2,1$ узлы решетки в основном состоянии даются формулой вида $x_{n}=f(n \rho)$, где эллиптич. функция $f(x)$, зависящая от констант $\rho, \lambda_{1}, \lambda_{2}$, имеет вещественный период 1. Спектр соответствуюшего оператора Шредингера в термодинамич. пределе $m, N \rightarrow \infty, m / N \rightarrow \rho$ состоит из трех разрешенных 30н



\[

\left[E_{1}, E_{2}\right],\left[E_{3},-E_{3}\right],\left[-E_{2},-E_{1}\right], E_{1}<E_{2}<E_{3}<0 .

\]



При $\rho<1 / 2$ заполнена нижняя зона, а две верхние пусты. При $\rho>1 / 2$ заполненными оказываются две нижние зоны. Во всех случаях уровень Ферми лежит в запрешенной зоне, то есть в основном состоянии система оказывается диэлектриком.



В интегрируемом случае основное состояние оказывается вырожденным и в системе появляется помимо обычных звуковых волн дополнительная ветвь спектра возбуждений. Зависимость энергии основного состояния $E(\rho)$ от плотности электронов является аналитической. При малых неинтегрируемых возмущениях основное состояние оказывается по-прежнему квазипериодич. функцией с двумя периодами 1 и р. Энергия основного состояния становится разрывной во всех рациональных точках $\rho=p / q$. Причем величина разрыва имеет порядок $\sim e^{- \text {const } q}$. Такой же характер имеет и энергия «пининга» - энергия зацепления дополнительной звуковой волны за узлы решетки.



Лит.: [1] П а й е рл с Р., Квантовая теория твердых тел, пер. с англ., М., 1956; [2] Б разо в ск и й С. А., Дзя лош и н ск и й И. Е., К р иче в е р И. М., «Журнал эксперим. и теоретич. физики», 1982, т. 82, в. 8; [3] Бе лок о л о с Е. Д., «Теоретич. и математич. физика», 1980 , т. 45 , в. 2 , с. $268-74$.



И. М. Кричевер.



ПАО ФОРМУЛА - см. Спектр турбулентности.



ПАPA СИЛ - система двух сил $\boldsymbol{P}$ и $\boldsymbol{P}^{\prime}$, действуюших на твердое тело, равных по абсолютной величине и направленных параллельно, но в противоположные стороны, то есть



\[

\boldsymbol{P}^{\prime}=-\boldsymbol{P}

\]



П. с. не имеет равнодействующей, то есть ее нельзя заменить (а следовательно, и уравновесить) одной силой. Расстояние $l$ между линиями действия сил пары называется п л е ч о м П. с. Действие, оказываемое П. с. на твердое тело, характеризуется ее моментом, к-рый изображается вектором $\boldsymbol{M}$, равным по абсолютной величине $\boldsymbol{P l}$ и направленным перпендикулярно к плоскости действия П. с. в сторону, откуда поворот, совершаемый П.с., виден происходящим против хода часовой стрелки. Основное свойство П. с.: действие, оказываемое ею на данное твердое тело, не изменяется, если П. с. переносить куда угодно в плоскости пары или в плоскости, ей параллельной, а также если изменять абсолютную величину сил пары и длину ее плеча, сохраняя неизменным момент П.с. Таким образом, момент пары можно считать приложенным в любой точке тела. Две П. с. с одинаковыми моментами, приложенные к одному и тому же телу, механически эквивалентны одна другой. Любая система П.с., приложенных к данному твердому телу, механически эквивалентна одной П. с. с моментом, равным сумме векторов-моментов этих П. с. Если сумма моментов нек-рой системы П.с. равна нулю, то эта система П. с. является уравновешенной.



C. M. Tapz.



ПАРАБОЛИЧЕСКИЕ КООРДИНАТЫ - сМ. Координаты. ПАРАБОЛИЧЕСКОТО ЦИЛИНДРА ФУНКЦИИ см. Вебера уравнение.



ПАРАБОЛИЧЕСКОЕ УРАВНЕНИЕ - см. Дифференциальное уравнение с частными производными второго порядка.



ПАРАЛЛЕЛОГРАММ ПЕРИОДОВ - см. Двоякопериодическая функция. ПАРАМЕТРИКС - см. Волновое уравнение; асимптотическое решение.



ПАРАМЕТРИЧЕСКИЙ МЕТОД - ОДИН ИЗ осНоВНЫХ Методов геометрической теории функций комплексного переменного, заключающийся в описании конформных отображений посредством однопараметрических семейств отображений специального вида.



Трудности построения эффективного вариационного исчисления на классах однолистных функций коренятся в существенной нелинейности объекта исследований. Способ их преодоления, предложенный К. Левнером (К. Löwner, 1923), состоит в следующем. Пусть $L-$ совокупность голоморфных и однолистных в $U=\{z:|z|<1\}$ функщий $\varphi$, удовлетворяющих условиям $\varphi(0)=0, \varphi^{\prime}(0)>0$ и $|\varphi(z)|<1$ при $z \in U$. Внутренние задачи теории однолистных функщий легко формулируются как экстремальные задачи на классе $L$. Заметив, что $L$ образует полугруппу относительно операции композиции, К. Левнер, по аналогии с теорией групп Ли, поставил и исследовал вопрос об ее инфинитезимальном описании.



Пусть $\left\{\varphi_{t}\right\}_{t} \geqslant 0$ - однопараметрич. семейство в $L$ такое, что $\varphi_{0}(z) \equiv z$ и $t \mapsto \varphi_{t}(z)$ дифференцируемо, локально равномерно в $U$ относительно $z$. Тогда производная



\[

\left.\frac{d \varphi_{t}(z)}{d t}\right|_{t=0}=v(z)

\]



представляет собой голоморфную в $U$ функцию и называется инфинитезимальным преобразованием полугруппы $L$. Непосредственно из Швариа леммы следует, что



\[

v(z)=-z p(z)

\]



где $p$ - голоморфная в $U$ функция с неотрицательной действительной частью. Вопрос об инфинитезимальном описании полугруппы $L$ состоит в возможности представления каждого $\varphi$ из $L$ в виде композиции инфинитезимальных преобразований, то есть в виде $\varphi(z)=w(T, z)$, где $T \geqslant 0$, а $w=w(t, z)-$ решение дифференциального уравнения



\[

\frac{d w}{d t}=-w P(w, t)

\]



с начальным условием $\left.w\right|_{t=0}=z$ и функцией $P(z, t)$, измеримой по $t$ на $[0, T]$, голоморфной по $z$ в $U$ и удовлетворяющей условиям $P(0, t)=1, \operatorname{Re} P(z, t)>0$ при $z \in U$. Уравнение (1) известно как ура в не ни е Л е в не ра. Положительное решение вопроса об инфинитезимальном описании полугруппы $L$ дает возможность применения средств анализа к исследованию задач теории однолистных функций. В частности, широкий круг экстремальных задач на $L$ приводится посредством (1) к задачам оптимального управления.



П. м. возник в связи с решением гипотезы Бибербаха для третьего коэффициента. В дальнейшем этот метод привел к решению широкого круга задач теории однолистных функций, среди к-рых проблема врашения и проблема коэффициентов.



Лит.: [1] Гол у з и Г Г. М., Геометрическая теория функций комплексного переменного, 2 изд., М., 1966.



В. В. Горяйнов.



ПАРАМЕТРИЧЕСКИЙ ОСЧИЛЛЯТОР - см. Осцилляmop.



ПАРАМЕТРИЧЕСКИЙ РЕЗОНАНС - явление потери устойчивости линейной гамильтоновой системой, зависящей от параметров, при периодическом изменении параметров во времени. Система, автономная и устойчивая при каждом фиксированном значении параметров, становится неустойчивой при периодич. изменении параметров подходящего периода и сколь угодно малой амплитуды. Для возникновения П. р. период изменения параметров должен быть таким, чтобы исходная автономная система, рассматриваемая как периодическая с этим периодом, имела два совпадающих мультипликатора (необходимо лежаших на единичной окружности ввиду устойчивости) и эти мультипликаторы были разного рода. Последнее означает, что при замене матрицы $A$ системы на матрицу $\lambda A$ и при изменении $\lambda$ мультипликаторы движутся по единичной окружности в разные





\end{document}